\paragraph{Put the cohort on an explicit calendar.}
The balances above specify a half-service observation but do not determine its
date or the sequence of transactions. For a reproducible dated version, supply
coverage from 1 January 2024 to 1 January 2025, exclusive, and observe the book
at the start of 2 July 2024. There are 183 elapsed days in this 366-day coverage
year, so uniform service gives exactly one half. These dates are synthetic
premises, not a chronology recovered from an insurer's disclosures.

Recognise the issued and held groups separately on 1 January and record their
premium receipts and payments then. At the July observation, first release
issued coverage and recognise incurred claims, then record the claim payment;
release held coverage, recognise the recovery measurement, and record the
recovery receipt. Together with the two premium movements, these are eight
ordered events. The group identifiers and the held-to-issued link accompany
every movement. Their order and opening balances are checked at each step, so
repeating a release cannot create another period of revenue.

This dated version gives the same cash 101, assets 119, liabilities 72 and
profit 27. Gross receipts are 135 and payments 54; their difference of 81
reconciles opening cash 20 to closing cash. The example supplies an operating
cash classification. Changing transferred risk from 1 to 1.5 increases only
the held recovery measurement and profit by 0.5; a separate non-performance
effect of 2 reduces them by 2. Neither change alters the recorded cash or nets
the two groups. The proposed company-book interface and complete interim report
are specified in Section~\ref{sec:insurance-paa-dated-book}. They retain unknown
actual accounting admission, complete notes and the unobserved remainder of
the year.

